\section{Control Architecture}
\label{sec:disc-control}

The \SI{1}{\mega\hertz} carrier is generated by the MCU carrier timer,
TIM1, while the MCU burst-enable pin PA10 controls transfer intervals.
ADC measurements set average demand without shifting the active carrier
away from the selected resonant region.

Circular DMA replaced blocking ADC conversion waits. This reduces
acquisition-related delay, which can otherwise admit an additional energy
packet before inhibition. The firmware attempts to align enable requests with carrier
boundaries to limit partial intervals. These measures control packet
timing; soft switching still requires matched commutation-current and
node-voltage measurements.

The half-bridge current input PA2 was initially misinterpreted. Its product
with secondary voltage cannot represent USB-C output power. In the baseline
firmware, secondary-voltage feedback PA1 closes the voltage loop, while
PA2 supports load classification and diagnostics.

Loaded demand is initialized from the stored average duty and trimmed by
voltage error. Voltage sag invokes additional duty; the upper threshold
suppresses transfer. Separate no-load caps and holdoff thresholds reduce
charging after detachment. Optional primary-voltage feed-forward caps
demand as the sensed link rises, while PA1 remains the feedback input.

DMA acquisition, carrier alignment, and loaded control were combined in
the final tests, so their individual effects were not isolated. A matched
comparison requires one feature to be varied while hardware, settings,
line condition, and load remain fixed.
